\section{有限维子空间逼近}
\label{sec:rkhs-finite-subspaces}

上一节从读数出发，目标函数未知，问题是在可行函数中选一个。本节换一个角度：目标函数$f\in H$已知，或者至少它与若干固定函数的内积可以算出，我们只允许用一个有限维子空间$V\subseteq H$中的函数来表示它。问题变成表示误差：$V$中离$f$最近的函数是谁，误差有多大，$V$逐步扩大后误差能否趋于零。

在Hilbert空间中，离$f$最近的元素就是正交投影$P_Vf$，误差为$\lVert f-P_Vf\rVert_H$（第~\ref{chap:function-inner-product-spaces}章）。RKHS的额外结构是，这个范数误差还能控制每一点的误差，而控制的系数只依赖于核与子空间，不依赖于$f$。

\begin{theorem}[子空间的再生核与逐点误差]
\label{thm:rkhs-subspace-error}
设$V$是RKHS $H$的闭子空间，$P_V$为到$V$的正交投影。则$V$带着$H$的内积是RKHS，其再生核为$K_V(x,y)=\langle P_VK_x,P_VK_y\rangle_H$；并且对一切$f\in H$与$x\in X$，
\begin{equation}
\lvert f(x)-(P_Vf)(x)\rvert\leq\sqrt{K(x,x)-K_V(x,x)}\ \lVert f-P_Vf\rVert_H .
\label{eq:rkhs-subspace-pointwise}
\end{equation}
\end{theorem}

\begin{proof}
对$g\in V$，$g(x)=\langle K_x,g\rangle=\langle P_VK_x,g\rangle$，所以$V$中$x$处的核截面是$P_VK_x$，再生核为式中所写。

令$r=f-P_Vf\in V^\perp$。$P_VK_x\in V$与$r$正交，所以
\[
r(x)=\langle K_x,r\rangle=\langle K_x-P_VK_x,r\rangle .
\]
用Cauchy--Schwarz不等式，并由勾股定理计算第一个因子：
\[
\lVert K_x-P_VK_x\rVert^2=\lVert K_x\rVert^2-\lVert P_VK_x\rVert^2=K(x,x)-K_V(x,x).
\]
这就是式~\eqref{eq:rkhs-subspace-pointwise}。
\end{proof}

式~\eqref{eq:rkhs-subspace-pointwise}右端的第一个因子是核截面$K_x$中不能由$V$表示的部分的长度，它只取决于核与子空间；第二个因子是$f$本身的表示误差。下面分别考察两类有限子空间：由核截面张成的子空间，以及由一般有限特征张成的子空间。

\subsection{核截面张成空间的逼近}
\label{sec:rkhs-section-span-approx}

取有限点组$Z=\{x_1,\ldots,x_n\}$，$V=M_Z=\operatorname{span}\{K_x:x\in Z\}$。由上一节，$P_{M_Z}f$就是$f$在$Z$上读数的最小范数插值，只需$f$的点值就能算出。这时式~\eqref{eq:rkhs-subspace-pointwise}中的第一个因子记为$P_Z(x)$，称为\term{核插值幂函数}（power function）。

$P_{M_Z}K_x$是$K_x$读数$\bm k_Z(x)=(K(x_1,x),\ldots,K(x_n,x))^{\top}$的最小范数插值，由定理~\ref{thm:rkhs-min-norm-interpolation}，
\begin{equation}
P_Z(x)^2=K(x,x)-\bm k_Z(x)^*\bm G^{\dagger}\bm k_Z(x).
\label{eq:rkhs-power-function}
\end{equation}
用矩阵的话说，这是$\bm K$关于观测下标的Schur补的对角元。它在每个观测点处为零，在观测点之间随间隔增大而增大。

\paragraph{点集决定的稠密性}

定理~\ref{thm:rkhs-kernel-properties}已经说明，全部核截面的组合在$H$中稠密，所以任何$f\in H$都能被有限个截面的组合任意逼近。实际中可选的位置往往受限，例如只能在一个点集$D\subseteq X$上放置截面。这时稠密性由下面的条件刻画。

\begin{theorem}[限定点集的核截面稠密性]
\label{thm:rkhs-dense-sections}
设$D\subseteq X$。则$\operatorname{span}\{K_x:x\in D\}$在$H$中稠密，当且仅当$H$中在$D$上处处为零的函数只有零函数。
\end{theorem}

\begin{proof}
由再生性质，$f$与每个$K_x$（$x\in D$）正交，当且仅当$f$在$D$上处处为零。一个子空间在Hilbert空间中稠密，当且仅当与它正交的元素只有零元。
\end{proof}

这个条件同时涉及核与点集。若$X$是度量空间且$K$连续，由再生性质与Cauchy--Schwarz不等式，
\[
\lvert f(x)-f(y)\rvert\leq\lVert f\rVert_H\lVert K_x-K_y\rVert_H ,
\]
其中核截面之差的长度可由核直接算出：
\[
\lVert K_x-K_y\rVert_H^2=K(x,x)+K(y,y)-2\operatorname{Re}K(x,y) .
\]
它随$y\to x$趋于零，所以$H$中的函数都连续。连续函数在稠密集上为零就处处为零，于是对连续核，任意稠密点集$D$上的截面都在$H$中稠密。

若核不连续，稠密的点集未必足够。取$X=[0,1]$上的离散核：$x=y$时$K(x,y)=1$，否则为$0$。这是“无限大的单位矩阵”，对应的RKHS是$\ell^2([0,1])$，即只在至多可数个点上非零且平方可和的函数。取$D$为$[0,1]$中的有理数，单点示性函数$\mathbf 1_{\{x_0\}}$（$x_0$为无理数）在$D$上为零而非零，所以$D$上的截面不稠密。核的连续性让读数在相邻位置之间传递信息，点集的稠密性才能发挥作用。

\begin{example}[最小值核的均匀节点]
\label{ex:rkhs-hmin-grid}
取步长$h=\tfrac1m$，在$H_{\min}$中取节点$Z_m=\{\tfrac1m,\tfrac2m,\ldots,1\}$。$P_{Z_m}f$是$f$的折线插值，其导数在每个小区间上等于$f'$在该区间上的平均。

对$q(x)=4x(1-x)$，$q'=4-8x$在每个长为$h$的区间上与其平均之差是斜率为$-8$、中点为零的线性函数，因此
\[
\lVert q-P_{Z_m}q\rVert_{H_{\min}}^2=m\int_{-h/2}^{h/2}(8u)^2\,\mathrm du=\frac{16h^2}{3}.
\]
$m=3$时为$\tfrac{16}{27}$，与例~\ref{ex:rkhs-hmin-interpolation}一致，并随$h\to0$趋于零。

幂函数也可以直接算出。设$x$落在小区间$[x_{k-1},x_k]$中，记$a=x-x_{k-1}$，$b=x_k-x$。$K_x-P_{Z_m}K_x$是$\min(\cdot,x)$减去它的折线插值，只在这个小区间上非零，斜率先为$b/h$、再为$-a/h$。把导数平方积分，
\[
P_{Z_m}(x)^2=a\cdot\frac{b^2}{h^2}+b\cdot\frac{a^2}{h^2}=\frac{ab}{h}.
\]
最大值$h/4$在区间中点取得。
\end{example}

\paragraph{RKHS范数误差与其他范数下的误差}

以上的误差都在$H$的范数下度量。若$\kappa^2=\sup_xK(x,x)<\infty$，由点评估界，
\[
\sup_x\lvert f(x)-g(x)\rvert\leq\kappa\lVert f-g\rVert_H ,
\]
所以$H$范数下的逼近蕴含一致逼近。反过来不成立。$H_{\min}$中$f_k(x)=\sin(k\pi x)/k$满足$\sup\lvert f_k\rvert\leq1/k$，而它的核范数不趋于零：
\[
\lVert f_k\rVert_{H_{\min}}^2=\int_0^1\pi^2\cos^2(k\pi x)\,\mathrm dx=\frac{\pi^2}2 .
\]
$H$范数还计入了函数的振荡，比一致范数更强。

另一个问题是用核截面去逼近$H$之外的函数，例如$[0,1]$上一个给定的连续函数$g$。这时误差只能在一致范数等较弱的范数下度量，能否任意逼近取决于$H$在相应函数空间中是否稠密，这是核本身的性质。仿射核的空间是二维的，在$C[-1,1]$中是闭子空间，不能一致逼近$\lvert x\rvert$。$H_{\min}$中的函数都在$0$处为零，只能一致逼近满足$g(0)=0$的连续函数。

高斯核的RKHS限制到紧集上以后，在该紧集的连续函数空间中按一致范数稠密，这类核称为\term{通用核}（universal kernel）。证明需要用到一致逼近的稠密性判据或测度的Fourier表示，篇幅较长，读者可查阅相关的核方法专著。它的含义是：只要节点选得足够多，高斯核截面的组合能在一致范数下逼近紧集上任意给定的连续函数。

所以核与点集各有分工。核决定$H$中有哪些函数、$H$能在多大范围内逼近外部函数，以及点值能否在相邻位置之间传递；点集决定所选截面能否充分利用$H$，即在$H$内部是否稠密。

\subsection{有限特征表示}
\label{sec:rkhs-finite-features}

核截面子空间的好处是坐标可以由点值读出，但它未必是表示给定函数类的最经济方式。另一种做法是在$H$中直接选定$m$个线性无关的函数$e_1,\ldots,e_m$作为特征，令$V=\operatorname{span}\{e_1,\ldots,e_m\}$，坐标预算就是$m$。

\paragraph{投影的计算}

记$\bm B$为特征在$H$内积下的Gram矩阵，元素为$\langle e_k,e_l\rangle_H$，它可逆。对$f\in H$，记$\bm b$为以$\langle e_k,f\rangle_H$为分量的向量。投影$P_Vf=\sum_la_le_l$的系数由正交条件确定：
\[
\langle e_k,f-P_Vf\rangle=0\qquad(k=1,\ldots,m),
\]
也就是$\bm B\bm a=\bm b$。误差由勾股定理得到：
\begin{equation}
\lVert f-P_Vf\rVert_H^2=\lVert f\rVert_H^2-\bm b^*\bm B^{-1}\bm b .
\label{eq:rkhs-feature-projection-error}
\end{equation}

对$V$的再生核，把$f=K_y$代入，$\bm b_k=\langle e_k,K_y\rangle=\overline{e_k(y)}$。记$\bm e(x)=(e_1(x),\ldots,e_m(x))^{\top}$，得到
\begin{equation}
K_V(x,y)=\bm e(x)^{\top}\bm B^{-1}\overline{\bm e(y)} .
\label{eq:rkhs-feature-subspace-kernel}
\end{equation}

与核截面相比，一般特征的坐标需要内积$\langle e_k,f\rangle_H$，它们一般不是点值读数。取$e_k=K_{x_k}$时，$\bm b_k=f(x_k)$，$\bm B$就是Gram矩阵$\bm G$，这时投影只需读数，这正是核截面的特殊之处。显式特征则可以利用对目标函数类的先验了解，用更少的坐标达到更小的误差。

\begin{example}[两个多项式特征]
\label{ex:rkhs-hmin-polynomial-features}
在$H_{\min}$中取$e_1=x$，$e_2=x^2$，它们都在$0$处为零。由$\langle f,g\rangle=\int_0^1f'g'$，
\[
\bm B=\begin{bmatrix}1&1\\1&\tfrac43\end{bmatrix},\qquad
\lVert a x+bx^2\rVert_{H_{\min}}^2=a^2+2ab+\tfrac43b^2 .
\]
$q(x)=4x-4x^2$本身属于$V$，两个坐标就给出零误差；例~\ref{ex:rkhs-hmin-interpolation}用三个核截面，误差却是$\sqrt{16/27}\approx0.77$。另一方面，这组特征只对二次多项式有利，换成振荡的目标函数，它未必比均匀节点的核截面更好。
\end{example}

\paragraph{有限特征核不是原核}

有限特征还可以用来定义一个新核。把例~\ref{ex:rkhs-hmin-polynomial-features}的特征排成向量$(x,x^2)$，按欧氏内积得到有限特征核$\widetilde K(x,y)=xy+x^2y^2$。由定理~\ref{thm:rkhs-feature-space}，$H_{\widetilde K}$与$V$含有同样的函数，范数却是$a^2+b^2$，与$H_{\min}$中的$a^2+2ab+\tfrac43b^2$不同。

$V$作为$H_{\min}$子空间的正确的核由式~\eqref{eq:rkhs-feature-subspace-kernel}给出。先求逆，再代入：
\[
\bm B^{-1}=\begin{bmatrix}4&-3\\-3&3\end{bmatrix},\qquad
K_V(x,y)=4xy-3xy^2-3x^2y+3x^2y^2 .
\]
它在对角线上满足$K_V(x,x)\leq K(x,x)=x$，例如$K_V(\tfrac12,\tfrac12)=\tfrac7{16}<\tfrac12$。由式~\eqref{eq:rkhs-subspace-pointwise}，任意$f\in H_{\min}$在$\tfrac12$处的投影误差不超过$\tfrac14\lVert f-P_Vf\rVert$。差$K-K_V$是$V^\perp$的再生核，只有$V=H$时才为零。

所以有限特征近似不应自动理解为原核的精确表示。特征的欧氏内积给出的是另一个空间的几何，只有按$H$的Gram矩阵加权，才得到原空间中子空间的核；即使如此，它也只代表原空间的一部分。
